% !TeX program = lualatex

\documentclass{article}
\usepackage[a4paper, total={5.5in, 10in}]{geometry}
\usepackage[utf8]{inputenc}
\usepackage{mathtools}
\usepackage[english]{babel}
\usepackage{amsthm}
\usepackage{calrsfs}
\usepackage{lmodern} %get scalable font
\usepackage{titling}
\usepackage{listings}
\usepackage{fancyhdr}
\usepackage{graphicx}
\usepackage[dvipsnames]{xcolor}
\usepackage{fontspec}
\usepackage{lstfiracode}
\usepackage[]{stmaryrd}

% Define command for show :=
\newcommand{\defeq}{\vcentcolon=}
% Set font family
\fontfamily{qpl}\selectfont 
% Calligraphic H
\DeclareMathAlphabet{\pazocal}{OMS}{zplm}{m}{n}
\newcommand{\Hb}{\pazocal{H}}
\newcommand{\Bb}{\pazocal{B}}

\newcommand*{\QEDB}{\null\nobreak\hfill\ensuremath{\square}}


\newtheorem{theorem}{Theorem}[section]
\newtheorem{corollary}{Corollary}[theorem]
\newtheorem{lemma}[theorem]{Lemma}
\newtheorem{definition*}[theorem]{Definition}

\title{Fifth hands-on: Bloom Filters \\[1ex] \large Algorithm Design (2021/2022)}
\author{Gabriele Pappalardo\\Email: g.pappalardo4@studenti.unipi.it\\Department of Computer Science}
\date{March 2022}

\pretitle{
    \begin{center}
    \LARGE
    \includegraphics[width=5cm]{/Users/gabryon/Projects/Papers/AD-21-22/assets/images/unipi.png}\\
}
\posttitle{
    \end{center}
}


\setmonofont[Scale=0.85,Contextuals=Alternate]{FiraCode}

\lstset{
    style=FiraCodeStyle,
    breakatwhitespace=false,         
    breaklines=true,                 
    captionpos=b,                    
    keepspaces=true,                
    showspaces=false,                
    showstringspaces=false,
    basewidth=0.6em,
    showtabs=false,                  
    tabsize=2,
    frame=single,
    basicstyle=\ttfamily, % Use \ttfamily for source code listings
    commentstyle=\color{ForestGreen}
}

\setcounter{section}{0}

\begin{document}

\maketitle

\section{Introduction}

The problem is composed in two parts:

\begin{enumerate}
    \item Consider the Bloom Filters where a single random universal hash random function $h : U \to [m]$ is employed for a set $S \subseteq U$ of keys, where $U$ is the universe of keys.
    Consider its binary array $B$ of $m$ bits. Suppose that $m \ge c|S|$, for some constant $c > 1$, and that both $c$ and $|S|$ are unknown to us.   
    Estimate the expected number of 1s in $B$ under a uniform choice at random of $h \in \mathcal{H}$. Is this related to $|S|$? Can we use it to estimate $|S|$? 
    
    \item Consider $B$ and its rank function: show how to use extra $O(m)$ bits to store a space-efficient data structure that returns, for any given $i$, the following answer in constant time: 
    $\textrm{rank}(i) = \#1s \in B[1..i]$
    
    \textit{Hint}: Easy to solve in extra $O(m \log m)$ bits. To get $O(m)$ bits, use prefix sums on B, and sample them. Use a lookup table for pieces of B between any two consecutive samples. 
\end{enumerate}

\section{Solution}

\subsection{Bloom Filter}

A Bloom Filter is a probabilistic data structure, invented in 1970 by Burton Bloom, that allows to check wheter an element $x$ belongs to a set $S$ without storing it.
A filter works with $k > 0$ hash functions $h_i: U \to [m]$ belonging to a universal hash family.

\subsection{Estimate expected number of bits set}

The first point of the hands on asks us to estimate the expected number of bits set in a Bloom Filter using only one hash function.
We start defining a new random indicator variable $X_i$ such that:
\begin{equation}
    X_i = \begin{cases}
        1 & \textrm{if} \, B_i = 1 \\
        0 & \textrm{otherwise}
    \end{cases}
\end{equation}

\noindent Therefore, we can be build a new random variable $Y = \sum_{i = 0}^{m - 1}X_i$ to estimate the number of bits set to one. We define $n = |S|$, as the unknown value to find. 

\begin{align*}
    E \llbracket Y \rrbracket = E \bigg\llbracket  \sum_{i = 0}^{m - 1}X_i \bigg\rrbracket = \sum_{i = 0}^{m - 1} E \llbracket X_i \rrbracket \\
    = \sum_{i = 0}^{m - 1} 1 - \bigg(1 - \frac{1}{m}\bigg)^n = (m - 1)\bigg(1 - \bigg(1 - \frac{1}{m}\bigg)^n\bigg) \\
    = m\bigg(1 - \bigg(1 - \frac{1}{m}\bigg)^n\bigg) \\
    = m - m\bigg(1 - \frac{1}{m}\bigg)^n \simeq m \bigg(1 - e^{-\frac{n}{m}}\bigg)
\end{align*}

\noindent Knowing that $\mu = E \llbracket Y \rrbracket \simeq m (1 - e^{-\frac{n}{m}})$, we can solve the equation to find the cardinality of the original set $S$, which is equal to $n$.

\begin{align*}
    \mu \simeq m (1 - e^{-\frac{n}{m}}) \\
    \iff \frac{\mu}{m} = 1 - e^{-\frac{n}{m}} \\
    \iff \frac{\mu}{m} - 1 = - e^{-\frac{n}{m}} \\
    \iff 1 - \frac{\mu}{m} = e^{-\frac{n}{m}} \\
    \iff \ln (1 - \frac{\mu}{m}) = -\frac{n}{m} \\
    \iff m \ln (1 - \frac{\mu}{m}) = - n \\
    \iff n = -m \ln (1 - \frac{\mu}{m})
\end{align*}

\noindent Thus, the number of elements in $S$ is $n =  -m \ln (1 - \frac{\mu}{m})$.

\subsection{Rank Computation}

The second point of the hands-on asks us to compute the rank of the bit array $B$, using only $O(m)$ bits of space. To reach the requested space complexity,
we start off by the baseline solution, the one gave by the hint.\\

\subsubsection{Baseline Solution}

\noindent Prefix sums allow us to answer the rank function in constant time, in fact, we can build a new array $P$, that contains in each position $i$ the number of ones
up to $i$.

\begin{equation}
    P_i = \sum_{j = 0}^{i} B_j
\end{equation}

\noindent Since the maximum prefix sum can be $|B| = m$ (when all the bits are set to $1$), we need $O(\log m)$ bits to store each sum, bringing us to use
$O(m \log m)$ space to store the entire array of prefix sums.\\

\subsubsection{Requested Solution}

\noindent To achieve the requested space $O(m)$ we sample the prefix sums array $P$, and we create a new lookup table $T$, described below.\\

\noindent First, we have to determine how many samples of the prefix sums we do need. We are going to use $\frac{2m}{\log m}$ samples of prefix sums. We will save
these samples in an array called $P$ of size $|P| = \frac{2m}{\log m}$.\\

\noindent Second, how do we build $T$? When we are asking for the rank of a position which has not a corresponding sample, 
we have to use the nearest sample to the position, plus the number of ones between the sample and the position itself.\\

\noindent To do so, we define $L = \frac{|B|}{\frac{2m}{\log m}} = \frac{\log m}{2}$ between two samples. We can split the array $B$ in $L$ parts.
Splitting the bit array leave us with $|B|/L$ portions, that we can use to index our lookup table $T$. 
How many strings of bits can be represented using $L$ bits? That is, how many rows $T$ will have? Exactly: $\#\textrm{rows} = 2^L$.

\begin{equation}
    \textrm{\#rows} = 2^L = 2^{\frac{\log m}{2}} = (2^{\log m})^{\frac{1}{2}} = \sqrt m
\end{equation}

\noindent We need $\sqrt m$ rows for $T$. How many columns will have the table $T$? $\# \textrm{columns} = L = \frac{\log m}{2}$. Therefore, our lookup table $T$ will have a size of $\sqrt m \frac{\log m}{2}$.
Assuming we have a bit array of size $m = 64$, then the lookup table will look like as shown below.

\begin{align}
     \begin{matrix}
         (0 & 0 & 0)_2 \\
         (0 & 0 & 1)_2 \\
         (0 & 1 & 0)_2 \\
         (0 & 1 & 1)_2 \\
         (1 & 0 & 0)_2 \\
         (1 & 0 & 1)_2 \\
         (1 & 1 & 0)_2 \\
         (1 & 1 & 1)_2
     \end{matrix} \quad
     T = 
     \begin{bmatrix}
        0 & 0 & 0 \\
        0 & 0 & 1 \\
        0 & 1 & 1 \\
        0 & 1 & 2 \\
        1 & 1 & 1 \\
        1 & 1 & 2 \\
        1 & 2 & 2 \\
        1 & 2 & 3
    \end{bmatrix} 
\end{align}

\noindent Our table will contain the partial sums of each of the possible portions as shown in Equation 4.

\noindent At the end, to compute the rank function we are going to need
both prefix sums samples and the lookup table $T$.

\subsubsection{Space Complexity}

We are keeping $\frac{2m}{\log m}$ samples of prefix sums, which takes up to $O(\log m)$ bits,
therefore $O(m)$.
Plus, for the lookup table $T$ we are using $\sqrt m$ rows and $\frac{\log m}{2}$ columns, and for each cell we
need $O(\log \log m)$ bits, thus $O(\sqrt{m}\log m (\log \log m))$ bits.

\begin{lstlisting}[language=C,caption=`Rank function implemented using a bit array of 64 elements.']
    uint64_t rank(uint64_t B, uint64_t j) {
        // Assuming we have the sample prefix sums array P
        // and the lookup table T.

        uint64_t m = 64; 
        uint64_t L = log2(m) / 2;
        
        // Row and Column indecies for the Lookup Table T
        uint64_t r = get_row(B, j);
        uint64_t c = j % L;

        return P[(j / L) - 1] + T[r][c];
    }

    uint64_t get_row(uint64_t B, uint64_t i) {
        // Shift and mask to get the row index
        return (B >> (0x40 - (0x3 * (i + 1)))) & 0x7;
    }
\end{lstlisting}

\noindent At the end we have $O(m) + O(\sqrt{m}\log m (\log \log m)) = O(m)$ total bits.

\end{document}
